\begin{textsample}{NEG Dim 5 – vem\_ed\_30 – Score 2.00 – vem\_ed\_30/t155.txt}  \label{ex:f5_neg_138}
Aos Leitores Osvaldo Succi Jr.
Coordenador dos PCIs Transformação digital na educação superior: esse foi o tema da International Digital Week 2025, organizada pela Université Côte D´Azur de 16 a 20 de junho.
Estive no evento e apresentei workshop sobre a revolução da Inteligência Artificial nos projetos COIL (Collaborative Online International Learning), compartilhei um workshop sobre Internacionalização Responsável com a professora Luciane Stallivieri (UFSC), participei de uma mesa-redonda e apresentei um pôster sobre o Projeto Colaborativo Internacional (PCI) desenvolvido entre Fatec Itatiba e Thomas More University of Applied Sciences (Bélgica).
Nesse projeto, premiado como “Melhor Pôster do Dia” em 19 de junho, os estudantes elaboraram um game educativo sobre preservação dos recursos hídricos.
Maria Claudia Nunes Delfino, responsável por Projetos Colaborativos Internacionais (PCIs) em inglês em nossa equipe de Apoio à Internacionalização do Ensino Superior na Divisão de Extensão e Pesquisa (Depes) da Coordenadoria Geral de Ensino Superior de Graduação (CGESG) do CPS, apresentou seis trabalhos na 13th International Corpus Linguistics Conference, que ocorreu de 30 de junho a 3 de julho em Birmingham, Reino Unido.
Nesta edição, conheça ainda o PCI desenvolvido entre Fatec Santo André e DUOC UC (Chile), sobre emissões de poluentes em veículos leves.
E o Artigo de Opinião sobre “IA nos Projetos Colaborativos Internacionais”, aprofundando um pouco da discussão que levei à International Digital Week.
Quer contribuir você também com artigo para VEm?
Acesse https://publicacoescesu.cps.sp.gov.br/VEm/about/submissions.
Dúvidas?
Escreva para cgesg.pci@cps.sp.gov.br.
Boa leitura!

% matched lemmas: 
\end{textsample}
